% E3 statistical strengthening snippet
\paragraph{Statistical tests of the window interpretation.}
To distinguish a stage-localised sensitive-window effect from a smooth early-to-late trend, I reanalysed E3 with three post hoc robustness tests: paired bootstrap comparisons between the E1/E2 window (steps 512--3000) and late checkpoints (steps 8000 and 143000), segmented-versus-monotone regression models, and uptake-controlled regressions. Step 0 and step 128 are retained as diagnostic baselines but excluded from the main model-comparison tests because uptake at these near-initial checkpoints is qualitatively different from trained checkpoints.
For normalised clean-retention margin, the window--late difference was 0.279 with a bootstrap 95\% interval [0.237, 0.318].
For normalised degradation-resistance AUC, the window--late difference was 0.203 with a bootstrap 95\% interval [0.191, 0.215].
These analyses use uptake-normalised outcomes, so the reported durability effect is not simply a comparison of raw post-injection scores. The central test is whether injection stage predicts the fraction of the injected signal that survives after matched continuation or adversarial overwrite.